\section{Introduction}
\label{sec:introduction}

\subsection{Context and motivation}

Operating room planning coordinates limited surgical capacity with clinical
teams, equipment and postoperative care. Decision horizon, patient pathway,
resources and performance criteria define distinct scheduling problems
\citep{cardoen2010,guerriero2011}. Improving room utilisation can shift pressure
to downstream units, which motivates integrated resource planning
\citep{rachuba2026}. Duration uncertainty adds a second difficulty: a schedule
built from estimates may execute differently, with delays, overtime or cases
unable to start \citep{denton2007}. These outcomes require separate metrics;
utilisation alone cannot summarise performance \citep{schouten2023}.

The Fil rouge AI and Healthcare project investigates this problem using
historical surgical activity. Its organisational motivation is to move from
calendar-driven booking towards capacity-aware admission and intervention
date recommendations \citep{group2026dataaudit}. The present experiments
address components of that ambition rather than a complete admission system.
One model assigns patients to multi-day operating sessions and penalises
session overload, bed demand and imbalance. A second assigns a daily workload
to rooms and evaluates minute-resolution execution in Mesa. It includes
turnover and temporary closures but excludes beds, staffing and speciality
compatibility. Their objective values are not directly comparable.

\subsection{Research question and positioning}

The research question is: \emph{with common information and evaluation rules,
how do operating room metaheuristics compare, and do improvements in predicted
schedule quality survive execution under duration deviations and temporary
room unavailability?}

Exact methods, metaheuristics and predictive--prescriptive pipelines are
already established in the literature \citep{fei2010,vieira2025,lex2025}.
Our contribution is a reproducible comparison with explicit representation,
information and resource assumptions. In the room experiments, every method
uses a shared decoder that fixes within-room order by increasing predicted
duration. Realised durations remain hidden during ordinary planning. Thus,
the study compares allocation search and execution policies within that
representation, not unrestricted surgical sequencing.

\subsection{Contributions and report organisation}

The report brings together three experimental contributions:
\begin{enumerate}
\item Five searches---simulated annealing, tabu search, a genetic algorithm,
      hybrid Tabu--Simulated Annealing and ant colony optimisation---compared
      with a greedy baseline and a small exactly enumerable room reference.
\item Static and closure-triggered reactive Mesa execution from shared
      initial schedules, plus a separate median/oracle duration comparison.
\item A length-of-stay (LOS) study linking learned estimates to session
      planning, with paired uncertainty scenarios and granularity sensitivity.
\end{enumerate}
The room campaign controls per-request fitness calls, whereas the LOS
campaign retains unequal source budgets. The hybrid search combines tabu
selection with annealing acceptance; simulation coupling instead evaluates
schedules during execution. Neither form of hybridisation implies clinical
validation or a universal best method.

The literature review and problem description establish scope.
Sections~\ref{sec:formulation}--\ref{sec:experimental-design} present the
formulations, methods and protocols; Sections~\ref{sec:results}
and~\ref{sec:discussion} report and interpret the evidence, followed by the
conclusion. The eleven-part structure keeps each experimental contribution
within these common sections.
